Las apps Ionic con Capacitor necesitan sustitutos creíbles de plugins nativos para probar en Node; si los mocks quedan dispersos, la cobertura engaña. Se propone un \textbf{módulo de simulación} para \tcode{@capacitor/geolocation}, \tcode{@capacitor/camera} y \tcode{@capacitor/preferences}, con perfiles de éxito, error y permisos, integrado en Jest o Jasmine mediante un bridge \tcode{Proxy}. Hay un \textbf{procedimiento para subir cobertura de líneas y ramas} y una \textbf{capa MCP} opcional para repetir mediciones \cite{mcp-spec2025}. La validación contrasta línea base y escenario con simulador en un proyecto Ionic de referencia \cite{nie2023gui_mobile_mapping,schafer2024testpilot,olsthoorn2024syntest_javascript}; no reemplaza pruebas en dispositivo. Diseño cuantitativo antes/después según el capítulo de metodología.
